\section{Related Work}
\label{sec:related}

\para{Probabilistic neural-symbolic inference.} DeepProbLog integrates neural predicates into ProbLog so neural outputs can serve as probabilistic facts inside symbolic programs \cite{manhaeve2018deepproblog}. It helped establish \mnist arithmetic as a benchmark for weakly supervised perception plus logic. A-NeSI keeps discrete symbolic semantics while using approximate inference networks to scale beyond the small cases where exact probabilistic inference is practical \cite{vankrieken2022anesi}. Our experiment follows this line in spirit, but it uses a deliberately small exact marginalization layer rather than a full probabilistic programming stack. This choice lets us isolate the core question of unseen output-label extrapolation.

\para{Declarative and differentiable logic.} Scallop provides a Datalog-like language for neurosymbolic programs with provenance semantics and differentiable reasoning \cite{li2023scallop}. Logic Tensor Networks ground first-order logic into differentiable objectives, and LTNtorch makes that approach available in PyTorch \cite{badreddine2022logic,carraro2024ltntorch}. These systems are broader than our implementation: they support many rules, relations, and reasoning patterns. We instead use one exact rule, addition, and test whether the symbolic computation alone changes behavior when the label support shifts.

\para{Relational and concept-level generalization.} Neural Logic Machines learn tensorized logical operations over objects and relations, targeting length and size generalization in relational domains \cite{dong2019neural}. The Neuro-Symbolic Concept Learner connects visual perception, semantic parsing, and symbolic program execution for visual question answering \cite{mao2019neuro}; later concept-intervention work studies whether symbolic concepts can make corrections more meaningful \cite{stammer2021right}. These papers address richer reasoning settings. Our study is narrower, but the narrowness is useful: digit grounding can be measured directly, and every OOD sum has an exact symbolic explanation.

\para{AddMNIST as a compositional probe.} \mnist arithmetic remains useful because it has a clean latent structure: the image labels are not the target labels, but the target follows from a known rule \cite{lecun1998gradient,manhaeve2018deepproblog}. Many \addmnist evaluations report in-distribution accuracy or use larger arithmetic variants. We focus on an explicit support gap: high-plus-high pairs are removed from training, so the learner never receives supervised examples for sums 14--18. This makes the evaluation less about memorizing arithmetic labels and more about composing learned digit predicates.
